\documentclass[11pt,a4paper]{article}
\usepackage[T1]{fontenc}
\usepackage{lmodern,amsmath,amssymb,amsthm,mathtools,booktabs,microtype,xcolor}
\usepackage[margin=26mm]{geometry}
\usepackage[colorlinks=true,linkcolor=blue!45!black,citecolor=blue!45!black,urlcolor=blue!45!black]{hyperref}
\newtheorem{theorem}{Theorem}[section]
\newtheorem{lemma}[theorem]{Lemma}
\newtheorem{proposition}[theorem]{Proposition}
\newtheorem{corollary}[theorem]{Corollary}
\theoremstyle{remark}\newtheorem{remark}[theorem]{Remark}
\newcommand{\Q}{\mathbb Q}\newcommand{\Z}{\mathbb Z}\newcommand{\HH}{\mathcal H}
\DeclareMathOperator{\Tr}{Tr}\DeclareMathOperator{\Nm}{N}
\setlength{\emergencystretch}{2em}
\title{Common local weights and quartics in 297 and 342 variables}
\author{}\date{9 November 2026}
\begin{document}\maketitle
\begin{abstract}
We specify two families of integer quartics by complete quadratic systems
and coefficient tables. The first has 297 total variables and 271 quadratic
rows; the second has 342 total variables and 314 quadratic rows. Their
local conditions of orders $1$, $10$ and $40$ use powers of one rational
weight. A seven-equation quadratic presentation of these powers replaces
separate rational weight computations. We prove the local extraction
statement, including the nonvanishing needed for cleared trace equations,
and simultaneous completeness at every finite place and at the real places.
Under the stated arithmetic premises, all sufficiently large members of
the 297-variable family define the integers in the rationals, and every
member of the 342-variable family with parameter $N\ge27464$ does so.
The latter gives an explicitly specified quartic with 4938 nonzero monomials
and coefficient height less than $2^{1791283}$.
\end{abstract}

\section{Statements and arithmetic scope}
All variable counts include the input coordinate $X$. Degree is total
degree. If $P$ has integer coefficients, $H(P)$ denotes their largest
absolute value. For $N\ge1$ put $A_N=80989^{2N+1}$. The fixed number
$A_N$ is a coefficient parameter, not a polynomial variable.

The coefficient tables \nolinkurl{quadratic-family297.json} and
\nolinkurl{quadratic-family342.json} specify primitive integer quadratic
rows $Q_{N,j}$ in respectively 297 and 342 coordinates. Define
\[
 P_N^{(297)}=\sum_{j=1}^{271}Q_{N,j}^2,
 \qquad P_N^{(342)}=\sum_{j=1}^{314}Q_{N,j}^2.
\]
The complete expanded coefficients are recorded in the files
\begin{center}
\nolinkurl{polynomial-family297.json}\\
\nolinkurl{polynomial-family342.json}.
\end{center}

\begin{theorem}\label{thm:main}
Under the arithmetic premises specified below, the following hold.
\begin{enumerate}
\item There exists $N_0\ge27464$ such that, for every $N\ge N_0$,
\[
 x\in\Z\quad\Longleftrightarrow\quad
 \exists\boldsymbol y\in\Q^{296},\ P_N^{(297)}(x,\boldsymbol y)=0.
\]
The proof does not determine a numerical sufficient value of $N_0$.
\item For every $N\ge27464$,
\[
 x\in\Z\quad\Longleftrightarrow\quad
 \exists\boldsymbol y\in\Q^{341},\ P_N^{(342)}(x,\boldsymbol y)=0.
\]
In particular $A_{27464}=80989^{54929}$ specifies a single such polynomial.
\item Both families have degree exactly four and content one, with
\[
 H(P_N)<2^{1520}80989^{4N+2}\qquad(N\ge1).
\]
For $N\ge27464$ the numbers of nonzero coordinate monomials are 4268
and 4938, respectively. The largest coefficient of $P_{27464}^{(342)}$
has exactly 1791283 binary digits and is the coefficient of $r^2$.
More precisely, both families have exact coefficient height
$80989^{4N+2}+1$ for $N\ge27464$.
\end{enumerate}
\end{theorem}

Here are the retained arithmetic premises, stated in full in
\cite[Section~2]{Source}. Let
\[
 K=\Q(\rho),\quad \rho^2=2,\quad d=3-\rho,\quad
 E:y^2=x^3-d^2x,\quad
 P=\left(\frac{57-28\rho}{8},\frac{-660+417\rho}{32}\right).
\]
We assume the specified group and height assertions
$2E(K)=\Z P$ and $h(x([n]P))\ge n^2-6$ for $n\ne0$;
the five-contact height inequality with
\[
 \HH=2^{110}210^{1800}\left(\prod_{p\le76176}p\right)^{19}+128
 <2^{2093448};
\]
the obstruction to negative odd valuations supplied by the normalized
pole equations outside $2,3,17$, together with their stated realization
for a positive input less than $17$ with square reduced denominator and
the required local squareclasses and valuations at $2,3,17$; and the
simultaneous prime-pattern and conjugate reduction-order assertions with
fixed finite congruences, real approximation, arbitrary fixed avoidance,
and the weighted separation property used by the source construction.
The selected contact primes are $1$ modulo $8\cdot3\cdot17$.
The weighted separation and fixed denominator dilation are justified
from the stated analytic estimates in \cite{Selection}. The present argument
uses no stronger prime-selection or pole-realization statement. In particular, every additional integer introduced in the
finite preparation is fixed before the contact primes are selected.

For precision, the five-contact radical of $t=u/v$ in lowest terms,
$v>0$, is the product of primes $p>76176$ having positive odd valuation
on one of $v,u,u-v,u-2401v,32u+343v$. Outside the four exceptional
finite contact values, the assumed inequality is
$h(t)\le\HH M_5(t)^{19}$. The normalized contact change below adds only
$\log3$ to the relevant height bound. The aggregate trace-height theorem
and effective elliptic sum-height theorem used in the two respective
families are proved in \cite{Aggregate,Sums}; their different effectiveness
properties account for the distinction in Theorem~\ref{thm:main}.
Finite coefficient verification alone proves none of the retained global
arithmetic premises.

\section{One weight for all required local orders}
Write $\sigma$ for conjugation on $K/\Q$. At an odd rational prime,
the valuations of $K$ are unramified and normalized by $v(p)=1$.
Let $a\in\Q^{\times2}$, let $t=a\lambda$, and put
\[
 \Theta=\{0,-1/2,-1,1/2399,-32/407\},\qquad
 F(t)=\prod_{\theta\in\Theta}(t-\theta),\qquad \delta=t-1.
\]
The retained quadratic system implies
\begin{equation}\label{eq:retained}
 F(t)e=(a\delta)^5,\qquad \delta^3q=e.
\end{equation}
These are identities among the represented quantities; the actual table
expresses their products by quadratic rows. On the domain $\delta\ne0$,
\begin{equation}\label{eq:weights}
 R=\frac{F(t)}{a^5\delta^6},\qquad
 q=\frac1{R\delta^4},\qquad q_m=q^{2m-1}\quad(m\ge1).
\end{equation}
The contact radical $M(t)$ is the product of primes $p>76176$ for which
$v_p(t-\theta)$ is positive and odd for some $\theta\in\Theta$.

\begin{lemma}\label{lem:domain}
Every solution of \eqref{eq:retained} has $F(t)\ne0$. If $p\mid M(t)$,
then $\delta\ne0$, $v_p(\delta)=0$, and $v_p(R)$ is odd. In particular
$q\ne0$. If $\delta=0$, then $M(t)=1$; if $q=0$, then $\delta=0$.
\end{lemma}
\begin{proof}
Every root of $F$ differs from $1$, so $F(t)=0$ would contradict the
first equation, since $a\ne0$. For $p>76176$, each nonzero difference
between distinct elements of $\Theta\cup\{1\}$ is a $p$-adic unit.
A contact prime thus divides just one factor of $F(t)$ to positive odd
order and makes $\delta$ a unit. Since $a$ is a rational square,
$v_p(F(t))-5v_p(a)$ is odd. All prime divisors of the reduced numerators
and denominators of $1-\theta$ are at most $76176$, proving $M(1)=1$.
Finally $q=0$ gives $e=0$, and the first equation forces $\delta=0$.
\end{proof}

\begin{lemma}\label{lem:min}
Let $v$ be an integer-valued discrete valuation, with $v(0)=+\infty$,
and suppose $v(q)$ is odd. For arbitrary positive integers $m_j$ and
field elements $D_j$, the nonzero summands of
\[
 T=1+\sum_{j=0}^{s-1}q^{(2m_j-1)2^j}D_j^{2^{j+1}}
\]
have distinct valuations. Hence $T\ne0$ and
\[
 v(T)=\min\bigl(0,2^j((2m_j-1)v(q)+2v(D_j)):0\le j<s\bigr)\le0.
\]
If $v(T)=0$, then $T$ is congruent to $1$ in the residue field.
\end{lemma}
\begin{proof}
Each finite nonconstant valuation is $2^j$ times an odd integer. Its
exact power of two divisor is therefore $2^j$, so it differs from zero
and from every valuation with a different index. The smallest valuation
is attained once. The ultrametric inequality is then an equality. If
that smallest value is zero, all the nonconstant terms have positive
valuation. A zero $D_j$ simply removes its term.
\end{proof}

\begin{proposition}\label{prop:extract}
Suppose $p>2$, $\nu=v_p(R)$ is odd, and $v_p(\delta)=0$.
Let $W=w_0+w_1\rho$ be a product of factors from Lemma~\ref{lem:min},
using $q$ from \eqref{eq:weights}. Suppose the rational quantities
$w_0,w_1,T_0,e,b$ satisfy
\begin{equation}\label{eq:trace}
 w_0^2-2w_0T_0-2w_1^2=0,\qquad e\delta b=T_0^2,
 \qquad e\delta=R^{-1}.
\end{equation}
Suppose that $b\ne0$ cannot have negative odd $p$-adic valuation.
Then $\nu>0$, and at one place $v\mid p$ every test in every factor
satisfies $v(D_j)\ge m_j$.
\end{proposition}
\begin{proof}
At every place above $p$, $v(q)=-\nu$ is odd. Every factor is nonzero
and has nonpositive valuation by Lemma~\ref{lem:min}. In particular
$W\ne0$. If $w_0=0$, the first cleared equation forces $w_1=0$,
a contradiction. It follows from that equation that
\[
 T_0=\frac{\Nm(W)}{2w_0}=\frac1{\Tr(W^{-1})}\ne0,
 \qquad b=\frac{R}{\Tr(W^{-1})^2}\ne0.
\]
Thus no inverse of $w_0$ or unproved nonvanishing assumption has been used.
Put $\ell=\max_{v\mid p}v(W)\le0$. The trace inequality gives
\[
 v_p(\Tr(W^{-1}))\ge-\ell,
 \qquad v_p(b)\le\nu+2\ell.
\]
The valuation of $b$ is odd. The obstruction therefore implies
$0<v_p(b)\le\nu+2\ell$, so $\nu>0$ and $\ell>-\nu/2$.
Choose a place attaining $\ell$. Each factor has valuation at least
$\ell$, because all factor valuations are nonpositive and sum to $\ell$.
Consequently every nonzero test term satisfies
\[
 2^j\bigl((1-2m_j)\nu+2v(D_j)\bigr)>-\nu/2.
\]
Thus
\[
 v(D_j)>\left(m_j-\frac12-\frac1{2^{j+2}}\right)\nu
 \ge m_j-\frac34.
\]
The last inequality uses $m_j\ge1$ and the positive odd integer $\nu\ge1$.
Integrality gives $v(D_j)\ge m_j$. The assertion is automatic if $D_j=0$.
\end{proof}

\section{Simultaneous completeness}
We now verify all places, including primes dividing new denominators.
This step is needed because equality of contact valuations alone would
not prove completeness.

Fix an integer input, the scalar parameter, the four integer root labels,
and the positive rational square scales $a_1,\ldots,a_r$. The retained
point preparation gives nonempty real and finite-place neighborhoods
on which all the rational test functions are defined. Choose smaller
neighborhoods with compact closure; at each of finitely many prepared
places their valuations have a fixed lower bound. Fix these neighborhoods
and the scalar preparation before choosing contact primes.

\begin{lemma}\label{lem:finite}
Let $S$ be a fixed finite set of rational primes containing the prepared
places and all required fixed denominators. There is a positive integer
$D$, fixed before prime selection, such that the following holds uniformly
for integers $u,w$ with $w\ne0$ and $u$ a prime outside $S$:
\[
 \lambda=\frac{u}{D^2w^2},\qquad t_i=a_i\lambda.
\]
For every $p\in S$, all the nonconstant terms of every modified factor
can be made arbitrarily small $p$-adically by the choice of $D$.
We may simultaneously require $D^2\theta/a_i\in\Z$ for every $i,\theta$.
\end{lemma}
\begin{proof}
First clear the finitely many indicated denominators. Increasing $v_p(D)$
then gives
\[
 v_p(\lambda)=-2v_p(D)-2v_p(w)\le-2v_p(D).
\]
Choose it so that, for each $i$,
\[
 v_p(t_i)<\min\bigl(0,\ v_p(\theta):\theta\in\Theta\setminus\{0\}\bigr).
\]
The ultrametric inequality gives
\[
 v_p(\delta_i)=v_p(t_i),\qquad v_p(F(t_i))=5v_p(t_i),
 \qquad v_p(e_i)=5v_p(a_i).
\]
Therefore
\begin{equation}\label{eq:smallq}
 v_p(q_i)=5v_p(a_i)-3v_p(t_i)
          =2v_p(a_i)-3v_p(\lambda)
          \ge2v_p(a_i)+6v_p(D).
\end{equation}
This lower bound tends to infinity with $v_p(D)$, uniformly in $u,w$.
For a test assigned order $m$ in position $j$, the valuation of its term
is
\[
 2^j\bigl((2m-1)v_p(q_i)+2v(D_j)\bigr).
\]
The tests have fixed valuation lower bounds in the chosen neighborhoods,
so every such expression tends to infinity uniformly. The same argument
at a place of $K$ above $p$ uses its normalized extension valuation;
ramification at two changes no inequality. There are finitely many
scales, places and tests, so one fixed integer $D$ works for all of them.
\end{proof}

The additional divisors of $D$ enter only the fixed avoidance modulus.
The coefficient list in the prime pattern is fixed after this choice:
\[
 u-\frac{\theta}{a_i}D^2w^2.
\]
Repeated coefficients are identified. The coefficient $0$, contributed
by $\theta=0$, makes $u$ itself one of the selected primes. The
simultaneous selection premise supplies these primes with avoidance of
the fixed data, the stated conjugate orders, and the fixed congruences.
Thus the hypothesis that $u$ is prime outside $S$ is part of the same
prime pattern, not a further independent restriction. A fixed enlargement
of $D$ changes neither the exponent in the weighted count nor the
permitted slowly decreasing separation parameter.

\begin{proposition}\label{prop:complete}
The retained scalar and prime preparation supplies simultaneous witnesses
for all the modified detectors. No additional lower bound on the scalar
parameter $N$ is needed for completeness.
\end{proposition}
\begin{proof}
Apply Lemma~\ref{lem:finite} after the fixed point neighborhoods and
scales have been chosen. At each detector, writing $v_*=D^2w^2$, the
exact factorization is
\begin{equation}\label{eq:Rproduct}
 R_i=\frac{\displaystyle\prod_{\theta\in\Theta}
                 (u-(\theta/a_i)v_*)}{v_*^5\delta_i^6}.
\end{equation}
Its denominator is a rational square. The selected primes avoid its
support and the fixed contact differences. Hence $v_p(R_i)=1$ at
its five selected contacts, and $v_p(R_i)$ is even elsewhere.

At a selected contact prime the designated primary place satisfies
$v(D_j)\ge m_j$ by the unchanged point and logarithm preparation.
Since $v(q_i)=-1$, every new term has positive valuation there.
Thus each factor is $1$ modulo the primary prime. At the opposite split
place Lemma~\ref{lem:min} says that every factor either has negative
valuation or is $1$ modulo the prime. Consequently the contributions
to $\Tr(W_i^{-1})$ are $1$ and either $0$ or $1$. The trace is a unit,
congruent to $1$ or $2$ modulo $p>76176$. The pole input
\[
 b_i=R_i\Tr(W_i^{-1})^{-2}
\]
has valuation one at each selected contact.

Outside those contacts, \eqref{eq:Rproduct} gives even valuation of
$R_i$; hence the valuation of $b_i$ is even, whether positive or negative.
In particular its reduced denominator is a square. This argument covers
all remaining primes, including divisors of the numerator or denominator
of $\delta_i$; no unit assertion about $\delta_i$ at such primes is needed.

At the finite prepared places, Lemma~\ref{lem:finite} makes every factor
arbitrarily close to $1$, and hence makes $\Tr(W_i^{-1})$ arbitrarily close
to $2$. Moreover the same calculation gives
\[
 v_p(R_i)=-v_p(t_i)-5v_p(a_i)
          =-v_p(\lambda)-6v_p(a_i),
\]
which can be made uniformly large by $D$. Thus $v_2(b_i)\ge3$,
$v_3(b_i)\ge1$ and $v_{17}(b_i)\ge2$ follow after a fixed choice of
$D$. At two choose the closeness of the trace to $2$ so that its valuation
is exactly one; at three and seventeen it is a unit. The local squareclasses
of $b_i$ equal those of $R_i$. Formula~\eqref{eq:Rproduct}, with the selected
primes congruent to $1$ modulo $8\cdot3\cdot17$, shows that these are
squares in $\Q_2,\Q_3,\Q_{17}$.

At both real embeddings, uniformly on the fixed compact point
neighborhoods, as $\lambda\to+\infty$,
\[
 R_i=\frac{1+O(\lambda^{-1})}{a_i^6\lambda}>0,\qquad
 q_i=a_i^2\lambda^{-3}(1+O(\lambda^{-1})).
\]
Thus each factor is positive for large $\lambda$ and tends to $1$.
A term of order $m$ in position $j$ decays as
$O(\lambda^{-3(2m-1)2^j})$. Hence
$\Tr(W_i^{-1})\to2$ and $b_i\to0^+$, so $0<b_i<17$.
The selection premise supplies arbitrarily large $\lambda$ after $D$
and all fixed congruences have been chosen. The pole-realization premise
therefore supplies fresh simultaneous witnesses. Equal scales may reuse
identical choices. All rational circuit quantities are defined on these
witnesses, which proves the proposition.
\end{proof}

\section{The quadratic presentation and exact counts}
The powers needed by the local factors have the presentation
\begin{equation}\label{eq:chain}
 \begin{gathered}
 q_2=q^2,\quad q_4=q_2^2,\quad cq_4=q,\quad q_8=q_4^2,\\
 cq_{19}=q_8^2,\quad q_{38}=q_{19}^2,\quad cq_{79}=q_{38}^2.
 \end{gathered}
\end{equation}
Every row has degree at most two. On $q\ne0$, the rows successively
force
\[
 q_2=q^2,\quad q_4=q^4,\quad c=q^{-3},\quad q_8=q^8,
 \quad q_{19}=q^{19},\quad q_{38}=q^{38},\quad q_{79}=q^{79}.
\]
There are no choices of roots or signs in this conclusion. At a consumed
contact prime $q\ne0$ by Lemma~\ref{lem:domain}. When $q=0$, that lemma
shows that the contact radical is empty; no division by $q$ is used in
soundness in this case. Completeness uses $\delta\ne0$ and fills all
seven rows by their displayed rational values.

For the point and logarithm notation of the source, put
\[
 A_i=f(U_i)\ (1\le i\le4),\quad B_i=\sigma f(V_i)\ (0\le i\le4),
 \quad C_i=L_i-L_i',\quad I=L_1-\beta,\quad D_H=f(H),
\]
\[
 G=L_3^2-L_2(L_2-2)L_3L_4+(L_2-2)^2L_4^2-1,
 \qquad \beta=51X^2+2504.
\]
Here $f=1/x$ extends by zero at the identity, $U_0=H$, $V_0=H-P$,
and the retained exact addition rows give $Q_i=U_i-V_i$. The alternate
ratios $L_i'$ already include conjugation. The nine tests $A_i,B_i$ have
order one; $C_1,I$ have order ten; and $D_H,C_2,C_3,C_4,G$ have order forty.
For an ordered list, its factor is
\[
 \mathcal T(D_0,\ldots,D_{s-1})
 =1+\sum_{j=0}^{s-1}q^{(2m(D_j)-1)2^j}D_j^{2^{j+1}}.
\]
The two-detector system uses
\[
 \begin{gathered}
 (D_H,C_1,B_0),\quad(C_2,I,B_2),\quad(C_3,A_2,B_3),\\
 (C_4,A_3,B_4),\quad(G,A_4),\quad(A_1,B_1),
 \end{gathered}
\]
and the three-detector system uses
\[
 \begin{gathered}
 (D_H,C_1,B_1,B_0),\quad(C_2,I,B_2),\quad(C_3,A_2,B_3),\\
 (C_4,A_3,B_4),\quad(G,A_4,A_1).
 \end{gathered}
\]
Each of the sixteen tests occurs exactly once. Their field powers and
all the factor and product rows are exactly those specified and verified
in \cite{Packing}. They remain quadratic; only the scalar weights change.

In the two-detector system the shared field powers use ten rational
coordinates, while the three-detector system uses eighteen. Each field
square is represented by
\[
 (u+v\rho)^2=(u^2+2v^2)+(2uv)\rho,
\]
so it takes two quadratic output rows. The needed fourth and eighth
powers of the $B_i$ are linear combinations of existing formal-logarithm
coordinates, as proved coefficient by coefficient in \cite{Packing}.

The exact table transformation removes the source coordinates
$353,354,356$ and their translates by $48$ for each detector. The
coordinate with source identifier $355$ now denotes $c$. The source
coordinates $359,369,360$ denote $q_{19},q_{38},q_{79}$ respectively.
The existing coordinates $352,361,362$ denote $q,q_2,q_4$. A coordinate
$q_8$ is already present in each of the three detectors. In the
two-detector system it is added with source identifier $850+i$,
$i=0,1$. All other rows and coordinates are retained literally.
Consequently the systems have the following sizes:
\begin{center}
\begin{tabular}{lrrr}
\toprule
Height scales&Total coordinates&Quadratic rows&Nonzero monomials\\\midrule
$1,S_0^2$&297&271&4268\\
$1,S_0^2,(S_0+S_1)^2$&342&314&4938\\\bottomrule
\end{tabular}
\end{center}
Here $S=\sum_{i=1}^4x(Q_i)=S_0+S_1\rho$. The monomial counts hold for
$N\ge27464$; they are statements about coefficients after substituting
$A_N$, not additional variables.


\subsection{Primitive normalization of the scalar contact product}
Let $z_j$ denote a coordinate by its source identifier. The actual
contact product is the affine expression
\[
 \begin{aligned}
 r={}&-103171881284-2103220824z_{248}-21050600z_{487}\\
     &-7629700z_{488}-45771136z_{489}.
 \end{aligned}
\]
We use $r$ itself as the coordinate in place of $z_{489}$. This is an
invertible rational affine change: its coefficient of $z_{489}$ is
nonzero, and every other coordinate, including the input, is fixed.
The source coordinate occurs only in its own defining row and the
scalar normalizer row. After substitution and primitive normalization,
the latter is exactly
\[
                  R_s+80989^2rR_s-A_Nr=0.
\]
Every equation remains quadratic. Reverse substitution reproduces both
source rows exactly up to a nonzero rational scalar, so their rational
solution sets correspond bijectively. Taking sums of squares preserves
the required conjunction over $\Q$.

In the resulting quartic the coefficient of $r^2$ is $A_N^2+1$.
Every other coefficient is at most linear in $A_N$, and its sum of
absolute parameter coefficients is less than $2^{1520}$. For
$N\ge27464$ one has $A_N>2^{1520}+1$, so all those coefficients have
absolute value less than $A_N^2$. Thus, for both families,
\[
                        H(P_N)=80989^{4N+2}+1
                        \qquad(N\ge27464).
\]
This is an exact height formula, proved by the complete coefficient
comparison; it does not assert an effective defining threshold for the
two-detector family.

\section{Projection proofs}
The point, scalar, formal-logarithm and height arguments now apply with
exactly their original local orders. We give the numerical deductions
explicitly. Write the nonzero integer root labels as $n,a,t,b$, and put
$B=\max(|n|,|a|,|t|,|b|)$ and $c=|a-1|$. The retained scalar construction
and its exceptional-point exclusions imply
\[
 c\ge80989^{L_N},\qquad L_N=\left\lfloor\frac{187(2N-1)}8\right\rfloor+1,
 \qquad h(\beta)\le\HH c^{19/24}.
\]
Its equations are unchanged. The common-place conclusion of
Proposition~\ref{prop:extract} supplies the orders required by the
formal-logarithm argument. Thus the fortieth powers of every contact
radical divide
\[
 \Delta=t^2-a(a-2)tb+(a-2)^2b^2-1.
\]
If $n-\beta\ne0$, the tenth powers divide its nonzero reduced numerator
$J$, which satisfies $|J|\le2B\exp(h(\beta))$. These are the precise
arithmetic comparison conclusions used below; no new interpretation of
an unguarded point or ratio is introduced.

For three scales, the proved effective height theorem gives
\[
 B^2\le2h(S_0)+2h(S_0+S_1)+82.
\]
If the $k$th powers of all three radicals divide a nonzero integer $J$,
the contact-height theorem bounds the heights of $\lambda,S_0^2\lambda$
and $(S_0+S_1)^2\lambda$ by
$T=(\HH+\log3)|J|^{19/k}$. Since $\lambda\ne0$, both $h(S_0)$ and
$h(S_0+S_1)$ are at most $T$. As $\HH>92$, we obtain
\begin{equation}\label{eq:effective}
 B^2\le4T+82\le5\HH|J|^{19/k}.
\end{equation}
For $N\ge27464$, the independently verified fixed-prime estimates give
\[
 |a|>(40\HH)^{10},\qquad \log(48\HH)<|a|/40,
 \qquad h(\beta)<|a|/40.
\]
If $\Delta\ne0$, then $|\Delta|\le14B^4$. Equation~\eqref{eq:effective}
with $k=40$ yields
$B^2<20\HH B^{19/10}$, contradicting $B\ge|a|$. Therefore $\Delta=0$.
The elementary integral Lucas descent gives $\log B>|a|-3$.
If $n-\beta\ne0$, equation~\eqref{eq:effective} with $k=10$ implies
\[
 \log B<10\log(20\HH)+19h(\beta)<29|a|/40,
\]
a contradiction. Thus $n=\beta\in\Z$.

For two scales, the aggregate trace-height theorem supplies constants
$M_1>0,M_2\ge1$ with
\[
 B^2\le M_1h(S_0)+M_2.
\]
The additive constant is not effective in that proof. The same
height calculation for $\lambda,S_0^2\lambda$ gives
\begin{equation}\label{eq:ineffective}
 B^2\le C|J|^{19/k}+M_2,\qquad C=M_1(\HH+\log3).
\end{equation}
If $\Delta\ne0$, this implies
$B^2\le C14^{19/40}B^{19/10}+M_2$, impossible for sufficiently large $B$.
The scalar lower bound ensures $B\ge c-1\to\infty$ as $N\to\infty$.
Hence $\Delta=0$ for all sufficiently large $N$. For $B^2\ge2M_2$,
if $n-\beta\ne0$, \eqref{eq:ineffective} gives
\[
 \log B\le10\log(2C)+19\log2+19h(\beta)
 \le10\log(2C)+19\log2+19\HH c^{19/24}=o(c),
\]
contrary to the Lucas lower bound $\log B>c-4$. Therefore $n=\beta$
also in the two-scale system for every sufficiently large $N$.

In either case, writing $X=u/v$ in lowest terms, $v>0$, the integrality
of $51X^2+2504$ implies $v^2\mid51$. Since $51$ is squarefree, $v=1$.
This proves soundness. Proposition~\ref{prop:complete} and the retained
root and scalar preparation prove completeness for every integer input.
Over $\Q$ a sum of squares vanishes exactly when every quadratic row
vanishes. These observations prove the projection assertions in
Theorem~\ref{thm:main}.

\section{Coefficient verification and reproduction}
The producer \nolinkurl{build_quartic_shared_weights.py} reads the
pinned complete source tables, performs the displayed row changes, and
expands every square by exact integer arithmetic. The independent checker
\nolinkurl{check_quartic_shared_weights.py} does not import the producer.
It reconstructs every changed row from \eqref{eq:chain}, compares every
retained row, and expands the entire squared sum using unordered pairs
of monomials. It also checks all coordinate incidences, the input
coordinate, degree four and content one. No coefficient is sampled.

Before fixed-parameter powers are collected, the quartic tables contain
4269 and 4939 terms. Write the coefficient of any coordinate monomial as
$c_0+c_1A+c_2A^2$. In both tables the maximum of
$|c_0|+|c_1|+|c_2|$ has 1520 binary digits. Hence
$H(P_N)<2^{1520}A_N^2$. Since $A_N>2^{1520}+1$ for $N\ge27464$, the
highest nonzero parameter term dominates all lower ones, proving the
stated ordinary support counts. A coefficient identically equal to one
proves content one for every parameter. The scalar normalization proves that the sole coefficient of degree
two in $A$ is the coefficient $A^2+1$ of $r^2$; every other coefficient
is at most linear in $A$. The same coefficient-sum bound therefore
proves that $H(P_N)=A_N^2+1$ for $N\ge27464$. Direct integer evaluation
at $A=80989^{54929}$ gives 1791283 binary digits for this largest coefficient.

The finite verification includes exact rational substitutions in the
new weight chain and 61824 bounded valuation checks. These supplement
Proposition~\ref{prop:extract}, whose proof covers every odd valuation;
they are not a substitute for that proposition or the arithmetic premises.
The two reproduction programs use only the Python standard library.

\begin{thebibliography}{9}
\bibitem{Selection}
\emph{Separation and denominator conditions in quadratic prime patterns},
9 November 2026.
Source: \nolinkurl{prime_selection_20261109_v12.tex}.

\bibitem{Packing}
\emph{Simultaneous valuation conditions and explicit quartics},
9 November 2026. Complete field-power identities, point and logarithm
normalizations, and independent coefficient verification.
Source: \nolinkurl{mixed_valuation_tests_20261109_v11.tex}.
\bibitem{Source}
\emph{A forced prime condition and an explicit quartic in 392 variables},
8 November 2026. Arithmetic input contract, fixed-prime certificate,
scalar exclusions, and Lucas descent.
Source: \nolinkurl{quartic392_20261108_v8.tex}.
\bibitem{Aggregate}
\emph{Aggregate trace heights of elliptic multiples}, 8 November 2026.
Source: \nolinkurl{aggregate_trace_height_20261108_v9.tex}.
\bibitem{Sums}
\emph{Effective heights of elliptic abscissa sums}, 8 November 2026.
Source: \nolinkurl{effective_elliptic_sums_20261108_v10.tex}.
\bibitem{ExplicitSource}
\emph{An explicit quartic in 351 variables}, 9 November 2026.
Complete quadratic source table with SHA-256
\nolinkurl{452d063b6d7da598471aa1ec0302b9e37611a44168610592931ec56aac241966}.
\bibitem{ExistenceSource}
\emph{A family of quartics in 301 variables}, 9 November 2026.
Complete quadratic source table with SHA-256
\nolinkurl{9917be941135778816ce25ecb0cbc872b4ed7c9f3691d79355e5e8352667a4fe}.
\bibitem{Glazer}
E.~Glazer, \emph{An explicit quartic defining the integers in the rationals},
Version~13, manuscript and coefficient data, commit
\nolinkurl{5bbd74f6a479776a332cbe036809d5dddae63720},
\url{https://github.com/elliotglazer/diophantine-z-q}.
\end{thebibliography}
\end{document}
